\chapter{Opis biznesowy systemu}
\label{cha:opisBiznesowy}

OSK-Manager jest aplikacją webową służącą do zarządzania ośrodkiem szkolenia kierowców. System obejmuje obsługę szkół jazdy, użytkowników, kursantów, instruktorów, kursów, lekcji, pojazdów, dostępności, harmonogramu oraz wybranych danych rozliczeniowych. Poniżej przedstawiono wymagania funkcjonalne i niefunkcjonalne wynikające z zakresu projektu.

\section{Wymagania funkcjonalne}
\label{sec:wymaganiaFunkcjonalne}

Wymagania funkcjonalne określają funkcje, które powinien udostępniać system OSK-Manager.

\begin{enumerate}
  \item \textbf{WF-01. Logowanie użytkownika.} System powinien umożliwiać zalogowanie użytkownika za pomocą adresu e-mail i hasła.
  \item \textbf{WF-02. Obsługa sesji.} System powinien umożliwiać odświeżenie sesji, sprawdzenie aktualnego użytkownika oraz wylogowanie.
  \item \textbf{WF-03. Rejestracja użytkowników.} System powinien umożliwiać uprawnionym użytkownikom rejestrowanie kursantów i instruktorów.
  \item \textbf{WF-04. Role użytkowników.} System powinien rozróżniać role: kursant, instruktor, manager i administrator.
  \item \textbf{WF-05. Zarządzanie profilami.} System powinien umożliwiać edycję danych profilu użytkownika, w tym telefonu, opisu i avatara.
  \item \textbf{WF-06. Zarządzanie OSK.} System powinien umożliwiać managerowi tworzenie, edycję, listowanie i usuwanie własnych ośrodków szkolenia kierowców.
  \item \textbf{WF-07. Przypisanie użytkowników do OSK.} System powinien umożliwiać przypisywanie kursantów i instruktorów do wybranych szkół jazdy.
  \item \textbf{WF-08. Zarządzanie kursantami.} System powinien umożliwiać przeglądanie listy kursantów, szczegółów kursanta, notatek oraz numeru PKK.
  \item \textbf{WF-09. Zarządzanie instruktorami.} System powinien umożliwiać dodawanie, edycję, usuwanie i przeglądanie instruktorów przypisanych do OSK.
  \item \textbf{WF-10. Kwalifikacje instruktorów.} System powinien przechowywać informację, do jakich typów kursów uprawniony jest instruktor.
  \item \textbf{WF-11. Zarządzanie pojazdami.} System powinien umożliwiać dodawanie, edycję, usuwanie, zmianę statusu i dodawanie zdjęć pojazdów.
  \item \textbf{WF-12. Typy kursów.} System powinien udostępniać słownik typów kursów, na przykład kategorii prawa jazdy.
  \item \textbf{WF-13. Zarządzanie kursami.} System powinien umożliwiać tworzenie i edycję kursów powiązanych z kursantem, OSK i typem kursu.
  \item \textbf{WF-14. Uczestnictwo w kursie.} System powinien umożliwiać przypisanie kursanta do kursu oraz zmianę statusu uczestnictwa.
  \item \textbf{WF-15. Zarządzanie lekcjami.} System powinien umożliwiać tworzenie, edycję i przeglądanie lekcji teoretycznych oraz praktycznych.
  \item \textbf{WF-16. Reguły lekcji praktycznej.} System powinien wymagać przypisania pojazdu do lekcji praktycznej.
  \item \textbf{WF-17. Harmonogram szkoły.} System powinien prezentować harmonogram zajęć w ujęciu szkoły, instruktora lub kursanta.
  \item \textbf{WF-18. Dostępność instruktora.} System powinien umożliwiać definiowanie domyślnych godzin pracy, wyjątków dziennych, blokad czasu i urlopów instruktora.
  \item \textbf{WF-19. Wyliczanie wolnych terminów.} System powinien obliczać wolne okna czasowe instruktora na podstawie dostępności, blokad, urlopów i zaplanowanych lekcji.
  \item \textbf{WF-20. Zapobieganie kolizjom.} System powinien blokować planowanie nakładających się lekcji dla tego samego instruktora lub pojazdu.
  \item \textbf{WF-21. Wydarzenia instruktorskie.} System powinien umożliwiać tworzenie wydarzeń typu jazda lub teoria oraz przypisywanie do nich kursantów.
  \item \textbf{WF-22. Obsługa płatności.} System powinien przechowywać informacje o planie płatności kursu i dokonanych wpłatach.
  \item \textbf{WF-23. Oceny lekcji.} System powinien umożliwiać ocenę zakończonej jazdy przez uprawnionego kursanta.
  \item \textbf{WF-24. Panel kursanta.} System powinien umożliwiać kursantowi podgląd własnych kursów i zaplanowanych lekcji.
  \item \textbf{WF-25. Panel instruktora.} System powinien umożliwiać instruktorowi podgląd własnego harmonogramu i dostępności.
\end{enumerate}

\section{Wymagania niefunkcjonalne}
\label{sec:wymaganiaNiefunkcjonalne}

Wymagania niefunkcjonalne określają oczekiwane cechy jakościowe systemu OSK-Manager.

\begin{enumerate}
  \item \textbf{WN-01. Bezpieczeństwo dostępu.} Dostęp do chronionych funkcji systemu powinien wymagać poprawnej sesji użytkownika.
  \item \textbf{WN-02. Kontrola uprawnień.} System powinien ograniczać operacje zgodnie z rolą użytkownika oraz jego powiązaniem z konkretnym OSK.
  \item \textbf{WN-03. Izolacja danych szkół.} Manager powinien zarządzać wyłącznie szkołami, których jest właścicielem, a instruktor i kursant powinni widzieć tylko powiązane dane.
  \item \textbf{WN-04. Spójność danych domenowych.} System powinien wymuszać powiązanie lekcji z kursem, kursem z OSK i typem kursu, a pojazdów oraz instruktorów z właściwą szkołą.
  \item \textbf{WN-05. Brak podwójnych rezerwacji.} System powinien zapobiegać konfliktom terminów dla instruktorów i pojazdów.
  \item \textbf{WN-06. Walidacja danych wejściowych.} Dane przekazywane do API powinny być walidowane po stronie backendu niezależnie od walidacji formularzy w interfejsie.
  \item \textbf{WN-07. Bezpieczne przechowywanie sesji.} Token odświeżający powinien być przechowywany w ciasteczku HTTP-only, a klucze serwisowe Supabase wyłącznie po stronie backendu.
  \item \textbf{WN-08. Utrzymywalność architektury.} Backend powinien zachowywać podział na trasy, kontrolery, serwisy, schematy walidacji i warstwę dostępu do danych.
  \item \textbf{WN-09. Modularność frontendu.} Frontend powinien dzielić logikę na strony, komponenty, kompozycje, typy i warstwę BFF.
  \item \textbf{WN-10. Responsywność.} Interfejs powinien być używalny na typowych ekranach komputerów i urządzeń mobilnych.
  \item \textbf{WN-11. Czytelność interfejsu.} System powinien prezentować dane w sposób zrozumiały dla pracownika biura, instruktora i kursanta.
  \item \textbf{WN-12. Obsługa błędów.} System powinien zwracać jednoznaczne komunikaty błędów dla sytuacji takich jak brak uprawnień, błędne dane, konflikt terminu lub brak zasobu.
  \item \textbf{WN-13. Testowalność.} Kluczowe reguły domenowe, zwłaszcza dostępność i konflikty harmonogramu, powinny być możliwe do testowania automatycznego.
  \item \textbf{WN-14. Skalowalność.} System powinien umożliwiać obsługę wielu szkół, instruktorów i kursantów bez zmiany podstawowego modelu danych.
  \item \textbf{WN-15. Zgodność czasowa.} Operacje harmonogramu powinny być interpretowane w lokalnym czasie Polski, zgodnie z założeniem Europe/Warsaw.
\end{enumerate}
